\begin{theorem}[Gauge-space uniformity dependency route]
\label{thm:gauge-space-uniformity-dependency-route}
Every accepted $\GaugeSpaceUp$ packet consumes the $\MetricUp$ distance rows, $\UniformSpaceUp$ entourage rows, and $\QuasiUniformSpaceUp$ directional rows before it exports finite gauge uniformity. In carrier notation, the point row $X$, gauge list $Q$, diagonal row $D$, refinement row $R$, symmetry or directional row $S$, transport row $H$, continuation row $C$, provenance row $P$, and local naming row $N$ must all be displayed before any uniform-completion consumer may use the generated gauge surface.
\end{theorem}
\begin{proof}
Project the carrier of \autoref{def:gauge-space-carrier}. The gauge rows $Q$ do not stand alone: the carrier description places them beside metric distance input, generated uniform entourage data, and the directional quasi-uniform boundary.

Apply \autoref{thm:gauge-space-entourage-generation}. The diagonal and refinement rows $D,R$ are the finite entourage-generation surface, while $S$ distinguishes the ordinary $\UniformSpaceUp$ route from the $\QuasiUniformSpaceUp$ directional route.

Finally use \autoref{thm:gauge-space-uniformspace-sibling-anchor}. It identifies the generated gauge rows as the sibling handoff to the uniform and quasi-uniform interfaces, and the structural rows $H,C,P,N$ preserve that handoff without adding a primitive entourage filter, selected completion, quotient equality, or ambient uniformity oracle.
\end{proof}
